%% kapa: exam=1 round=30 year=2019 period=1 form=A subject=civil_law subject_ko=민법 seq=23 number=23
%% source: Q-Net 감정평가사 기출문제 (cst00309) artlSeq=5205205 file=extracted/1008/5205205/2019년 제30회 감정평가사 1차시험 문제지/2019년 제30회 감정평가사 1차시험 문제지/2019년도 제30회 감정평가사 1차 1교시 A형.hwp
%% conversion: kordoc 4.24.51 -> tools/gen1.py (auto); review=proofread-sonnet; flags=none
\begin{question}{23}
부동산의 물권변동을 위해 등기가 필요한 것은? (다툼이 있으면 판례에 따름)
\begin{choices}
  \item 건물의 신축에 의한 소유권의 취득
  \item 상속에 의한 토지 소유권의 취득
  \item 피담보채권의 소멸에 의한 저당권의 소멸
  \item 관습법에 따른 법정지상권의 취득
  \item 점유취득시효에 의한 토지 소유권의 취득
\end{choices}
\end{question}
\answer{5}
